\documentclass[11pt]{letter}
\usepackage[T1]{fontenc}
\usepackage[utf8]{inputenc}
\usepackage{lmodern}
\usepackage{geometry}
\geometry{margin=1in}

\signature{Haobo Ma and Wenlin Zhang}
\address{Wenlin Zhang (Corresponding Author)\\National University of Singapore\\Singapore\\e1327962@u.nus.edu}
\date{\today}

\begin{document}
\begin{letter}{Editorial Board\\Qualitative Theory of Dynamical Systems}

\opening{Dear Editors,}

We submit our manuscript entitled

\medskip
\centerline{\textbf{Tilt Dynamics of Cylinder Information and the Parry Measure}}
\centerline{\textbf{on the Golden-Mean Shift}}
\medskip

\noindent for consideration in \textit{Qualitative Theory of Dynamical Systems}.

The paper studies the fluctuation structure of cylinder information $\tau_n(x)=-\log\mu([x_0^{n-1}])$ for stationary Markov laws on the golden-mean shift, with emphasis on the zero-jitter condition $\mathrm{Var}(\tau_n)=o(n)$. The main contributions are:

\begin{itemize}
\item A full-parameter information normal form decomposing $\tau_n - nh(p)$ into an occupation-number bulk term and an endpoint correction, with the bulk coefficient vanishing precisely at the Parry point.

\item Tilt closure: exponential tilts of cylinder information preserve the one-step Markov family, and the induced dynamics on the parameter space linearize globally via the coordinate $L(p)=\log((1-p)/p^2)$.

\item A free-energy composition law $F_{q_p(t)}(s)=F_p(s+t-st)-(1-s)F_p(t)$, yielding a variance transfer formula and characterizing the Parry measure through constancy of its R\'enyi entropy spectrum.

\item Extensions to general mixing shifts of finite type: a Parry-relative fluctuation normal form, Parry rigidity in the full Markov simplex via Perron--Frobenius, and universality of the quadratic coefficient $\sigma_\tau^2 \sim 2(\log\lambda - h)$ near the Parry point.
\end{itemize}

We believe the manuscript fits the scope of QTDS because it combines explicit symbolic computations on a classical shift of finite type with rigidity results from thermodynamic formalism, and establishes structural phenomena (the composition law, the universal coefficient~2) that extend beyond the golden-mean setting.

The manuscript is original, has not been published previously, and is not under consideration elsewhere. All authors have approved the submission.

\closing{Sincerely,}

\end{letter}
\end{document}
